\documentclass[11pt]{article}

\usepackage{amsmath,amssymb,amsthm,mathtools}
\usepackage[margin=1in]{geometry}
\usepackage[hidelinks]{hyperref}

\newtheorem{lemma}{Lemma}

\title{A direct level-one AGT check on the 4-punctured sphere}
\author{}
\date{}

\begin{document}
\maketitle

\section*{Setup}

We take $\mathbb{P}^1$ with punctures at $(0,q,1,\infty)$ and external Liouville
momenta $(\beta_0,\alpha_0,\alpha_1,\beta_2)$. Write
\[
Q=\epsilon_1+\epsilon_2,\qquad
h(\alpha)=\frac{\alpha(Q-\alpha)}{\epsilon_1\epsilon_2},
\]
and set
\[
h_1=h(\alpha_0),\quad h_2=h(\beta_0),\quad h_3=h(\alpha_1),\quad h_4=h(\beta_2),
\quad h_{\mathrm{int}}=h\!\left(a+\tfrac{Q}{2}\right).
\]
Note $\epsilon_1\epsilon_2\,h_{\mathrm{int}}=\frac{Q^2}{4}-a^2$.

\section*{The two sides at instanton order one}

\paragraph{Virasoro side.}
At level~$1$ the only descendant is $L_{-1}|h_{\mathrm{int}}\rangle$; the sewing
formula gives
\begin{equation}
\label{eq:vir}
c_1^{\mathrm{Vir}}
=\frac{(h_{\mathrm{int}}+h_1-h_2)(h_{\mathrm{int}}+h_3-h_4)}{2h_{\mathrm{int}}}.
\end{equation}

\paragraph{Gauge-theory side.}
For SU$(2)$ $N_f=4$, the $k=1$ instanton moduli space has two torus-fixed
points, $([1],\varnothing)$ and $(\varnothing,[1])$. Evaluating the vector and
fundamental localization factors gives
\[
Z_1
=-\frac{\prod_{f=1}^{4}(a+\mu_f)}{\epsilon_1\epsilon_2\,(2a)(2a+Q)}
-\frac{\prod_{f=1}^{4}(-a+\mu_f)}{\epsilon_1\epsilon_2\,(2a)(2a-Q)},
\]
and the reduced block is $(1-q)^{-\nu}Z_{\mathrm{inst}}(q)$ with
$\nu=2\alpha_0(Q-\alpha_1)/(\epsilon_1\epsilon_2)$, so
\begin{equation}
\label{eq:agt}
c_1^{\mathrm{AGT}}=Z_1+\nu.
\end{equation}

\paragraph{AGT dictionary for the masses.}
\begin{equation}
\label{eq:dict}
\mu_1=-\tfrac{Q}{2}+\alpha_0+\beta_0,\quad
\mu_2=\tfrac{Q}{2}+\alpha_0-\beta_0,\quad
\mu_3=\tfrac{3Q}{2}-\alpha_1-\beta_2,\quad
\mu_4=\tfrac{Q}{2}-\alpha_1+\beta_2.
\end{equation}

\medskip

\noindent
The claim is $c_1^{\mathrm{Vir}}=c_1^{\mathrm{AGT}}$ under~\eqref{eq:dict}.
This is not trivial: \eqref{eq:vir} is a single rational function, while
\eqref{eq:agt} is a sum of three terms. The point of the check is to see how
the three-term structure collapses.

\section*{Key factorizations}

Set
\[
A(x):=(x+\mu_1)(x+\mu_2),\qquad B(x):=(x+\mu_3)(x+\mu_4),
\]
so $\prod_f(\pm a+\mu_f)=A(\pm a)\,B(\pm a)$.
Using~\eqref{eq:dict}, $\mu_1+\mu_2=2\alpha_0$ and
$\mu_1\mu_2=\alpha_0^2-(\beta_0-Q/2)^2$, so completing the square gives
\[
A(x)=(x+\alpha_0)^2-(\beta_0-\tfrac{Q}{2})^2,
\qquad
B(x)=(x+Q-\alpha_1)^2-(\beta_2-\tfrac{Q}{2})^2.
\]

Introduce the abbreviations
\[
P:=\epsilon_1\epsilon_2(h_{\mathrm{int}}+h_1-h_2),\qquad
R:=\epsilon_1\epsilon_2(h_{\mathrm{int}}+h_3-h_4),
\]
\[
s:=2a+Q,\qquad t:=2a-Q,\qquad
\sigma:=\alpha_0,\qquad \tau:=Q-\alpha_1.
\]
Note $s+t=4a$, $st=4a^2-Q^2=-2\epsilon_1\epsilon_2\,(2h_{\mathrm{int}})$, and
$\nu=2\sigma\tau/(\epsilon_1\epsilon_2)$.

\begin{lemma}
\label{lem:linear}
Under~\eqref{eq:dict},
\begin{align*}
A(a)  &= \sigma\,s - P, & A(-a) &= -\sigma\,t - P,\\
B(a)  &= \tau\,s - R,   & B(-a) &= -\tau\,t - R.
\end{align*}
\end{lemma}

\begin{proof}
A direct expansion:
\[
P=\epsilon_1\epsilon_2 h_{\mathrm{int}}+\alpha_0(Q-\alpha_0)-\beta_0(Q-\beta_0)
 =\tfrac{Q^2}{4}-a^2+\alpha_0(Q-\alpha_0)-\beta_0(Q-\beta_0),
\]
hence
\[
\sigma s - P
=\alpha_0(2a+Q)-\tfrac{Q^2}{4}+a^2-\alpha_0(Q-\alpha_0)+\beta_0(Q-\beta_0)
=a^2+2\alpha_0 a+\alpha_0^2-(\beta_0-\tfrac{Q}{2})^2
=A(a).
\]
Sending $a\to -a$ gives $A(-a)=-\sigma t-P$ (note $P$ is even in $a$).
The $B$ identities are the same computation with $(\alpha_0,\beta_0)\to(Q-\alpha_1,\beta_2)$
and $(h_1,h_2)\to(h_3,h_4)$.
\end{proof}

\section*{Collapse of the three-term sum}

With Lemma~\ref{lem:linear},
\[
A(a)B(a)=(\sigma s-P)(\tau s-R)=\sigma\tau s^2-(\sigma R+\tau P)s+PR,
\]
\[
A(-a)B(-a)=(-\sigma t-P)(-\tau t-R)=\sigma\tau t^2+(\sigma R+\tau P)t+PR.
\]
Combine $Z_1$ over the common denominator $\epsilon_1\epsilon_2(2a)\,s\,t$:
\[
Z_1=-\frac{A(a)B(a)\,t+A(-a)B(-a)\,s}{\epsilon_1\epsilon_2\,(2a)\,s\,t}.
\]
The numerator telescopes:
\begin{align*}
A(a)B(a)\,t+A(-a)B(-a)\,s
&= \sigma\tau\,st(s+t)
  +(\sigma R+\tau P)\bigl(-st+ts\bigr)
  +PR\,(s+t)\\
&= (s+t)\bigl(\sigma\tau\,st+PR\bigr)
 = 4a\bigl(\sigma\tau\,st+PR\bigr).
\end{align*}
Therefore
\[
Z_1
=-\frac{4a\bigl(\sigma\tau\,st+PR\bigr)}{\epsilon_1\epsilon_2\,(2a)\,st}
=-\frac{2\sigma\tau}{\epsilon_1\epsilon_2}-\frac{2PR}{\epsilon_1\epsilon_2\,st}
=-\nu-\frac{2PR}{\epsilon_1\epsilon_2\,(4a^2-Q^2)}.
\]
Adding $\nu$ kills the first term exactly:
\begin{equation}
\label{eq:agt-simplified}
c_1^{\mathrm{AGT}}=Z_1+\nu=-\frac{2PR}{\epsilon_1\epsilon_2\,(4a^2-Q^2)}.
\end{equation}

\section*{Comparison with the Virasoro side}

Since $\epsilon_1\epsilon_2\,(2h_{\mathrm{int}})=\tfrac{Q^2}{2}-2a^2=-\tfrac12(4a^2-Q^2)$,
\[
c_1^{\mathrm{Vir}}
=\frac{(h_{\mathrm{int}}+h_1-h_2)(h_{\mathrm{int}}+h_3-h_4)}{2h_{\mathrm{int}}}
=\frac{PR/(\epsilon_1\epsilon_2)^2}{2h_{\mathrm{int}}}
=\frac{PR}{\epsilon_1\epsilon_2}\cdot\frac{1}{\epsilon_1\epsilon_2(2h_{\mathrm{int}})}
=-\frac{2PR}{\epsilon_1\epsilon_2\,(4a^2-Q^2)}.
\]
This matches~\eqref{eq:agt-simplified}:
\[
c_1^{\mathrm{Vir}}=c_1^{\mathrm{AGT}}.
\]

\section*{Remarks}

\begin{enumerate}
\item The $U(1)$ factor $\nu$ is precisely what is needed to cancel the
$\sigma\tau\,st=\alpha_0(Q-\alpha_1)\,(4a^2-Q^2)$ piece in $Z_1$. This piece is
$a$-independent (after division by $4a^2-Q^2$), reflecting the fact that $\nu$
does not depend on the internal momentum $a$.
\item The identity $A(a)=\sigma s-P$ is simply a repackaging of the
3-point coefficient into ``propagator pole'' $\times$ ``external momentum'' plus
the Virasoro building block $\epsilon_1\epsilon_2(h_{\mathrm{int}}+h_1-h_2)$.
A parallel statement holds for the right vertex $(h_3,h_4)$, with
$\tau=Q-\alpha_1$ playing the role of $\sigma=\alpha_0$; the asymmetry
$\alpha_0\leftrightarrow Q-\alpha_1$ reflects the two different mass
assignments in~\eqref{eq:dict}.
\item The numerical identity has been double-checked symbolically against the
Nekrasov/Shapovalov series produced by \texttt{agt.py} (both sides agree as
rational functions of $\epsilon_1,\epsilon_2,a,\alpha_0,\alpha_1,\beta_0,\beta_2$).
\end{enumerate}

\end{document}
